%=============================================================================
%
%                   I-TYPE (INSTRUCTIONAL) PUBLICATION MODULE
%
%=============================================================================
%
% References:
%   - MIL-STD-38784C (General Style and Format Requirements for Technical
%     Manuals), PDF output requirements
%   - MARCORSYSCOM I-Type template
%
% One module serves MI (Modification), SI (Supply), TI (Technical) and LI
% (Lubrication) Instructions. They share a skeleton and differ only in which
% optional paragraphs appear, so the type is a label here rather than a fork.
%
% An I-Type is a directive, not correspondence: it has no From/To, it is
% identified by a PCN rather than an SSIC, and it is authenticated rather than
% signed off. Hence the empty addressing block below.
%
% Phase 1 scope is the frame -- title block, CUI markings and the
% authentication block. The numbered paragraph model, the nine materiel tables
% and the WARNING/CAUTION/NOTE environments follow; see
% docs/TECHPUB_I_TYPE.md.
%
%=============================================================================


%=============================================================================
%                             NO LETTERHEAD
%=============================================================================
% main.tex calls \printLetterhead unconditionally, which is right for
% correspondence: a letter is issued by a command and says so at the top. A
% technical publication is not -- it is identified by its nomenclature, short
% title and PCN, and its cover carries the Marine Corps seal at a fixed 2 x 2
% inches rather than a unit's address block.
%
% Overridden here rather than gated in main.tex so the twenty correspondence
% formats keep the behaviour they were written against. The cover page
% replaces this with the seal; see docs/TECHPUB_I_TYPE.md.

\renewcommand{\printLetterhead}{}


%=============================================================================
%                            TITLE BLOCK
%=============================================================================
% Long Title is centred and all caps per the MARCORSYSCOM template. The four
% line cap and the no-acronyms rule are author-facing, so they are enforced in
% the editor rather than silently truncated here.

\newcommand{\printDateAndTitle}{%
    \hfill \DocumentDate\par
    \vspace{12pt}%
    \begin{center}
        \MakeUppercase{\SubjectLine}
    \end{center}
    \par\vspace{12pt}%
}


%=============================================================================
%                        CONTROL MARKINGS BLOCK
%=============================================================================
% No From/To: an I-Type is issued to the fleet, not addressed to anyone. What
% would be the addressing block carries the control markings instead, matching
% where every other format puts them so the shared CUI macros keep working.

\newcommand{\printAddressBlock}{%
    \ifCUIEnabled
        \noindent
        Controlled by: \CUIControlledBy\\
        CUI Category: \CUICategory\\
        Limited Dissemination Control: \CUIDissemination\\
        \printPOCLine%
        \par\vspace{12pt}%
        \noindent\CUIDistStatement
        \par\vspace{12pt}%
    \fi
    %
    \ifClassifiedEnabled
        \noindent
        Classified by: \ClassifiedBy\\
        Derived from: \DerivedFrom\\
        Reason: \ClassificationReason\\
        Declassify on: \DeclassifyOn\\
        \printPOCLine
        \par\vspace{12pt}%
    \fi
}


%=============================================================================
%                          AUTHENTICATION BLOCK
%=============================================================================
% The template mandates the spacing exactly: two blank lines before OFFICIAL,
% four between OFFICIAL and the printed name, and two between the organisation
% title and DISTRIBUTION. At 12pt a blank line is 12pt, so the vspaces below
% are those counts made literal -- this is the sort of rule an author is asked
% to maintain by hand in Word, and the reason to generate the page at all.

\newcommand{\printSignature}{%
    \par\vspace{24pt}%      two blank lines
    \noindent OFFICIAL
    \par\vspace{48pt}%      four blank lines
    \noindent
    \ifNotEmptyElse{\SignatoryAbbrev}{\MakeUppercase{\SignatoryAbbrev}\\}{\ifNotEmpty{\SignatoryName}{\MakeUppercase{\SignatoryName}\\}}%
    % Signing authority then controlling office, per the template's
    % OFFICIAL block -- the reverse of the rank-then-title order the
    % correspondence formats use.
    \optionalLine{\SignatoryTitle}%
    \optionalLine{\SignatoryRank}%
    \par\vspace{24pt}%      two blank lines
    \noindent DISTRIBUTION: EDO
}


%=============================================================================
%                              PAGE STYLE
%=============================================================================
% Page numbers are centred in the footer per the MARCORSYSCOM template. Note
% this differs from MIL-STD-38784C PDF output (bold, lower right); the
% customer's own template is treated as authoritative -- see
% docs/TECHPUB_I_TYPE.md.

\fancypagestyle{firstpage}{%
    \fancyhf{}%
    \fancyhead[C]{\placeClassificationMarkings}%
    \fancyfoot[C]{}%
    \renewcommand{\headrulewidth}{0pt}%
    \renewcommand{\footrulewidth}{0pt}%
}

\fancypagestyle{documentpage}{%
    \fancyhf{}%
    \fancyhead[C]{\placeClassificationMarkings}%
    \fancyfoot[C]{\thepage}%
    \renewcommand{\headrulewidth}{0pt}%
    \renewcommand{\footrulewidth}{0pt}%
}
